\documentclass[11pt,a4paper]{article}
\usepackage[margin=22mm]{geometry}
\usepackage{amsmath,amssymb}
\usepackage{fontspec}
\usepackage{polyglossia}
\setdefaultlanguage{russian}
\setmainfont{DejaVu Serif}
\usepackage[unicode,hidelinks]{hyperref}

\allowdisplaybreaks[1]
\raggedright
\setlength{\parskip}{0.35em}
\newcommand{\dd}{\mathop{}\!\mathrm{d}}
\newcommand{\avg}[1]{\langle #1\rangle}
\newcommand{\ket}[1]{\lvert #1\rangle}
\newcommand{\bra}[1]{\langle #1\rvert}
\newcommand{\braket}[2]{\langle #1\mid #2\rangle}
\DeclareMathOperator{\Var}{Var}

\newcommand{\solutiontitle}[4]{%
  \hypersetup{pdftitle={Квантовая механика 1. Решения ДЗ №#1}}%
  \begin{center}
    {\small Квантовая механика 1 \enspace\textperiodcentered\enspace Семестр 1\par}
    \vspace{1ex}
    {\LARGE\bfseries Решения домашнего задания №#1\par}
    \vspace{0.5ex}
    {\large #2\par}
  \end{center}
  \noindent{\small Условия приведены в файле «#3», #4.}\par\medskip
}

\begin{document}
\solutiontitle{2}{Моменты измерений и функции операторов}{ДЗ2026(1)\_02.pdf}{с.~1--2}

\section{Собственные векторы и замена базиса}

Для \(A=\begin{pmatrix}1&2\\3&4\end{pmatrix}\) характеристическое
уравнение имеет вид \((1-\lambda)(4-\lambda)-6=\lambda^2-5\lambda-2=0\).
Значит,
\[
 \lambda_\pm=\frac{5\pm\sqrt{33}}2,\qquad
 v_\pm=\begin{pmatrix}2\\ \lambda_\pm-1\end{pmatrix}.
\]
Первая строка системы обнуляется при подстановке \(v_\pm\).
Вторая обнуляется по характеристическому уравнению. Два различных
корня дают два одномерных собственных подпространства.
Столбцы матрицы \(S=\begin{pmatrix}0&-1\\1&0\end{pmatrix}\)
суть координаты новых базисных векторов в старом базисе. Если
\(v=S v'\), то \(Av=ASv'=S A'v'\), откуда
\[
 A'=S^{-1}AS=\begin{pmatrix}4&-3\\-2&1\end{pmatrix}.
\]

\section{Сумма результатов измерений}

Пусть \(Q_j\) имеют общее конечное среднее \(\mu\) и конечную
дисперсию \(\sigma^2\). Линейность ожидания даёт
\(\avg{\tilde Q}=\sum_{j=1}^N\avg{Q_j}=N\mu\).
Раскрытие квадрата суммы центрированных результатов даёт дисперсию
\[
 \Var(\tilde Q)=N\sigma^2+
     2\sum_{i<j}\avg{(Q_i-\mu)(Q_j-\mu)}.
\]
Так что \(\Var(\tilde Q)=N\Var(Q)\) для независимых бросков или,
в более общем случае, при нулевой сумме попарных ковариаций. В условии формула для
дисперсии напечатана как повтор формулы для среднего. При зависимых
измерениях ковариации нужно учитывать, например при \(Q_j=Q_1\) для
всех \(j\) дисперсия суммы равна \(N^2\sigma^2\).

\section{Распределение Лоренца}

Пусть \(b=|a|>0\). Поскольку
\(\int_\mathbb R \dd x/(x^2+b^2)=[\arctan(x/b)/b]_{-\infty}^{\infty}
=\pi/b\), имеем \(C=b/\pi\).
На положительной полуоси
\[
 \int_0^R x\rho(x)\,\dd x=\frac{b}{2\pi}
                 \log\!\left(\frac{R^2+b^2}{b^2}\right)\longrightarrow +\infty.
\]
На отрицательной полуоси интеграл расходится к \(-\infty\).
Симметричное главное значение равно нулю, а математическое ожидание
не существует. Интегранд \(x^2\rho(x)\) стремится к \(b/\pi\) при
\(|x|\to\infty\), поэтому второй момент бесконечен и дисперсия не
определена. При \(a=0\) нормировка плотности невозможна из-за
расходимости при \(x=0\).

\section{Дисперсия по известным амплитудам}

По ортонормированности \(\|\Psi\|^2=|C|^2(1/6+1/2+1/3)=|C|^2\).
Берём \(C=1\). Вероятности значений \(1,2,3\) равны \(1/6,1/2,1/3\).
Поэтому
\[
 \avg{A}=\frac16+\frac22+\frac33=\frac{13}{6},\qquad
 \avg{A^2}=\frac16+\frac42+\frac93=\frac{31}{6},\qquad
 \Var(A)=\frac{31}{6}-\frac{169}{36}=\frac{17}{36}.
\]

\section{Вероятности по известным моментам}

Пусть \(p_j\) есть вероятность результата \(j\). Из
\(\avg{A^2}=\Var(A)+\avg{A}^2=9/20+25/4=67/10\) следует система
\[
 p_1+p_2+p_3=1,\qquad p_1+2p_2+3p_3=5/2,\qquad
 p_1+4p_2+9p_3=67/10.
\]
Вычитая первое уравнение из двух последних, получаем
\(p_2+2p_3=3/2,\ 3p_2+8p_3=57/10\).
Исключив \(p_2\), находим \(2p_3=57/10-9/2=6/5\).
Искомый набор \((p_1,p_2,p_3)=(1/10,3/10,3/5)\) положителен
и в сумме даёт единицу.

\section{Среднее, дисперсия и спектр матрицы}

Нормированный столбец \(\psi=(1,-i,i)^{\mathsf T}/\sqrt3\).
Умножив на заданную матрицу, получаем
\(K\psi=(1-i,1,i)^{\mathsf T}/\sqrt3\). Эрмитовость \(K\) даёт
\[
 \avg{K}=\psi^\dagger K\psi=\frac{(1-i)+i+1}{3}=\frac23,\qquad
 \avg{K^2}=\|K\psi\|^2=\frac43,\qquad \Var(K)=\frac89.
\]
Оператор переводит \(e_1\) в \(\sqrt2\,v\) и \(v\) в \(\sqrt2\,e_1\),
где \(v=(0,1,i)^{\mathsf T}/\sqrt2\). На ортогональном дополнении
\(u_0=(0,i,1)^{\mathsf T}/\sqrt2\) имеем \(Ku_0=0\).
Остальные собственные векторы \(u_\pm=(e_1\pm v)/\sqrt2\)
отвечают \(\pm\sqrt2\). Проекции состояния дают
\[
 \braket{u_0}{\psi}=\frac{-1+i}{\sqrt6},\qquad
 \braket{u_\pm}{\psi}=\frac{1}{\sqrt6}
                     \pm \frac{1-i}{2\sqrt3}.
\]
Вероятности \(p_0=1/3,\ p_\pm=1/3\pm 1/(3\sqrt2)\) положительны,
следовательно возможны все три значения.

\section{Экспоненты, обратный оператор и ядра (7*)}

\textbf{(a)} Для чётности \(P f(x)=f(-x)\) имеем \(P^2=I\).
Чётные степени в ряде экспоненты дают \(I\cos\alpha\), а нечётные
дают \(i P\sin\alpha\), поэтому
\[
 (e^{i\alpha P}f)(x)=\cos\alpha\,f(x)+i\sin\alpha\,f(-x).
\]
Ряд для \(e^{\alpha\partial_x}f\) есть ряд Тейлора \(f(x+\alpha)\)
там, где он сходится к функции. Для общего \(L^2\) унитарная группа
определена через Фурье-преобразование. Оно переводит \(\partial_x\)
в умножение на \(ik\), а множитель \(e^{i\alpha k}\) в обратном
преобразовании заменяет \(x\) на \(x+\alpha\). Значит
\((e^{\alpha\partial_x}f)(x)=f(x+\alpha)\).

\textbf{(b)} Предположим, что \(A^{-1}\) существует. Из
\((A-\lambda B)\sum_{n\ge0}\lambda^n C_n=I\) следуют
\(AC_0=I,\ AC_n=BC_{n-1}\) при \(n\ge1\).
Итак, \(C_n=(A^{-1}B)^n A^{-1}\) и
\[
 (A-\lambda B)^{-1}
 =A^{-1}+\lambda A^{-1}BA^{-1}
      +\lambda^2 A^{-1}BA^{-1}BA^{-1}+\cdots .
\]
При ограниченном \(A^{-1}B\) и \(|\lambda|\|A^{-1}B\|<1\)
ряд сходится по норме. Умножение геометрического ряда слева или
справа на \(I-\lambda A^{-1}B\) даёт \(I\). Вместе с
\(A-\lambda B=A(I-\lambda A^{-1}B)\) это проверяет обе стороны
обращения для матриц или ограниченных операторов с ограниченным
обратным \(A^{-1}\). В общем случае рекурсия задаёт формальный ряд
на общей области произведений.

\textbf{(c)} Для \(F(t)=e^{tA}Be^{-tA}\) дифференцирование с
сохранением порядка даёт \(F'(t)=e^{tA}[A,B]e^{-tA}\).
Если \([A,B]=cI\), то \(F'(t)=cI\) и \(F(0)=B\), значит
\(F(1)=B+cI\). Вычисление верно для сходящихся операторных
произведений или как тождество формальных рядов.

\textbf{(d)} По определению сопряжения и по перестановке
интегралов на пробных функциях
\begin{align*}
 \braket{f}{Lg}
  &=\iint f(x)^* L(x,x')g(x')\,\dd x\,\dd x'\\
  &=\int \left[\int L(x,x')^* f(x)\,\dd x\right]^*g(x')\,\dd x'
    =\braket{L^\dagger f}{g}.
\end{align*}
Отсюда \(L^\dagger(x,x')=L(x',x)^*\).
Свойство \(\int\delta(x-x')g(x')\dd x'=g(x)\) и производная по
\(x\) дают \(p_x(x,x')=-i\hbar\,\partial_x\delta(x-x')\).
Аналогично \(\int\delta(x+x')g(x')\dd x'=g(-x)\), значит
\(P(x,x')=\delta(x+x')\).

\end{document}
